\documentclass[11pt]{article}

\usepackage[margin=1in]{geometry}
\usepackage[T1]{fontenc}
\usepackage{lmodern}
\usepackage{microtype}
\usepackage{amsmath,amssymb}
\usepackage[dvipsnames]{xcolor}
\usepackage{enumitem}
\usepackage{needspace}
\usepackage[most]{tcolorbox}
\usepackage[hidelinks]{hyperref}

\hypersetup{
  pdftitle={Homework 2, MATH 327, Fall 2026},
  pdfauthor={Manish Patnaik},
  pdfsubject={MATH 327 homework problems}
}

\definecolor{courseblue}{HTML}{1F4E79}
\definecolor{lightblue}{HTML}{EAF2F8}
\definecolor{boxborder}{HTML}{B8CBDC}
\definecolor{darkgray}{HTML}{333333}

\setlength{\parindent}{0pt}
\setlength{\parskip}{0.55em}
\setlist[enumerate]{leftmargin=2em, itemsep=0.25em, topsep=0.3em}

\newtcolorbox{exercise}[1]{
  enhanced,
  colback=white,
  colframe=boxborder,
  colbacktitle=lightblue,
  coltitle=courseblue,
  fonttitle=\bfseries,
  title={Exercise #1},
  boxrule=0.6pt,
  arc=1mm,
  left=2.5mm,
  right=2.5mm,
  top=1.5mm,
  bottom=1.5mm,
  before skip=0.8em,
  after skip=0.8em
}

\newcommand{\chapterheading}[1]{%
  \par\vspace{0.5em}
  {\color{courseblue}\Large\bfseries #1}\par\vspace{0.35em}
  {\color{boxborder}\hrule height 0.6pt}\vspace{0.35em}
}

\newcommand{\sectionheading}[1]{%
  \Needspace{8\baselineskip}
  \vspace{0.25em}
  {\color{darkgray}\large\bfseries #1}\par
}

\begin{document}

\begin{center}
  \colorbox{lightblue}{%
    \begin{minipage}{0.92\textwidth}
      \vspace{0.6em}
      \begin{tabular*}{\textwidth}{@{\extracolsep{\fill}}lr}
        {\color{courseblue}\bfseries\LARGE Homework 2} &
        {\color{courseblue}\bfseries\LARGE MATH 327}\\[0.4em]
        \multicolumn{2}{l}{\color{darkgray}\bfseries Fall 2026}\\[0.35em]
        \multicolumn{2}{l}{\color{darkgray}Covers Lectures 8--13}\\[0.2em]
        \multicolumn{2}{l}{\color{darkgray}Suggested completion date: October 5}
      \end{tabular*}
      \vspace{0.5em}
    \end{minipage}}
\end{center}

\vspace{0.5em}

Unless otherwise indicated, the following exercises are from Michael Artin's
\emph{Algebra}, second edition.

\chapterheading{Chapter 2: Groups}

\sectionheading{Section 8: Cosets}

\begin{exercise}{8.1}
Let $H$ be the cyclic subgroup of the alternating group $A_4$ generated
by the permutation $(123)$. Exhibit the left and the right cosets of $H$
explicitly.
\end{exercise}

\begin{exercise}{8.4}
Does a group of order $35$ contain an element of order $5$? Of order $7$?
\end{exercise}

\begin{exercise}{8.10}
Prove that every subgroup of index $2$ is a normal subgroup, and show by
example that a subgroup of index $3$ need not be normal.
\end{exercise}

\sectionheading{Section 9: Modular Arithmetic}

\begin{exercise}{9.6}
Prove the Chinese Remainder Theorem: Let $a,b,u,v$ be integers, and assume
that the greatest common divisor of $a$ and $b$ is $1$. Then there is an
integer $x$ such that
\[
  x\equiv u\pmod a
  \qquad\text{and}\qquad
  x\equiv v\pmod b.
\]
\textit{Hint.} Do the case $u=0$ and $v=1$ first.
\end{exercise}

\newpage
\sectionheading{Section 10: The Correspondence Theorem}

\begin{exercise}{10.5}
With reference to the homomorphism $\varphi\colon S_4\to S_3$ described
in Example 2.5.13, determine the six subgroups of $S_4$ that contain $K$.

\textit{Recall.} The homomorphism $\varphi$ comes from permuting the three
partitions of $\{1,2,3,4\}$ into two pairs:
\[
\begin{aligned}
  \Pi_1&=\bigl\{\{1,2\},\{3,4\}\bigr\},&
  \Pi_2&=\bigl\{\{1,3\},\{2,4\}\bigr\},&
  \Pi_3&=\bigl\{\{1,4\},\{2,3\}\bigr\}.
\end{aligned}
\]
Here $K=\ker\varphi=\{1,(12)(34),(13)(24),(14)(23)\}$.
\end{exercise}

\sectionheading{Section 11: Product Groups}

\begin{exercise}{11.6}
Let $G$ be a group that contains normal subgroups of orders $3$ and $5$,
respectively. Prove that $G$ contains an element of order $15$.
\end{exercise}

\begin{exercise}{11.9}
Let $H$ and $K$ be subgroups of a group $G$. Prove that the product set
$HK$ is a subgroup of $G$ if and only if $HK=KH$.
\end{exercise}

\sectionheading{Section 12: Quotient Groups}

\begin{exercise}{12.4}
Let $H=\{\pm1,\pm i\}$ be the subgroup of $G=\mathbb{C}^{\times}$ of
fourth roots of unity. Describe the cosets of $H$ in $G$ explicitly.
Is $G/H$ isomorphic to $G$?
\end{exercise}

\begin{exercise}{12.5}
Let $G$ be the group of upper triangular real matrices
$\begin{bmatrix}a&b\\0&d\end{bmatrix}$ with $a$ and $d$ different from
zero. For each of the following subsets, determine whether or not $S$ is
a subgroup, and whether or not $S$ is a normal subgroup. If $S$ is a
normal subgroup, identify the quotient group $G/S$.
\begin{enumerate}[label=\textbf{(\roman*)}]
  \item $S$ is the subset defined by $b=0$.
  \item $S$ is the subset defined by $d=1$.
  \item $S$ is the subset defined by $a=d$.
\end{enumerate}
\end{exercise}

\newpage
\sectionheading{Miscellaneous Problems}

\begin{exercise}{M.12}
\textit{Postage stamp problem.} Let $a$ and $b$ be positive, relatively
prime integers.
\begin{enumerate}[label=\textbf{(\alph*)}]
  \item Prove that every sufficiently large positive integer $n$ can be
        obtained as $ra+sb$, where $r$ and $s$ are positive integers.
  \item Determine the largest integer that is not of this form.
\end{enumerate}
\end{exercise}

\sectionheading{Supplementary Problem: Semidirect Products (Optional)}

\begin{exercise}{S.1}
Let $G_1$ and $G_2$ be groups, and let
$\rho\colon G_2\to\operatorname{Aut}(G_1)$ be a homomorphism, where
$\operatorname{Aut}(G_1)$ is the group of automorphisms of $G_1$ under
composition. Write $\rho_{g_2}=\rho(g_2)$. Write out explicitly what it
means for $\rho$ to be a homomorphism.

The \emph{semidirect product} $G_1\rtimes_\rho G_2$ has underlying set
$G_1\times G_2$ and multiplication
\[
  (g_1,g_2)(g_1',g_2')
  =\bigl(g_1\rho_{g_2}(g_1'),\,g_2g_2'\bigr).
\]
\begin{enumerate}[label=\textbf{(\alph*)}]
  \item Verify that this law defines a group structure on
        $G_1\rtimes_\rho G_2$. What is the identity? What is the inverse
        of $(g_1,g_2)$?
  \item Identify $G_1$ and $G_2$ with the subgroups
        $G_1\times\{1\}$ and $\{1\}\times G_2$, respectively.
        Show that $G_1$ is normal in $G_1\rtimes_\rho G_2$.
        Is $G_2$ also normal?
  \item What choice of $\rho$ recovers the direct product?
  \item Show that the dihedral group $D_{2n}$ of order $2n$
        (the symmetry group of a regular $n$-gon, $n\geq3$) is
        isomorphic to a semidirect product. What are $G_1$, $G_2$,
        and $\rho$?
  \item If $N\trianglelefteq G$ and $H\leq G$ satisfy
        $N\cap H=\{1\}$ and $NH=G$, we say that $G$ is an
        \emph{internal semidirect product} of $N$ and $H$.
        What is the corresponding homomorphism
        $\rho\colon H\to\operatorname{Aut}(N)$? Explain how it gives
        an isomorphism $N\rtimes_\rho H\cong G$.
\end{enumerate}
\end{exercise}

\vfill
{\small\color{darkgray}Compiled \today.}

\end{document}
